Deep networks can fit arbitrary labels, including random ones, which makes memorisation of mislabelled training examples a central worry when labels are noisy~\cite{zhang2016understanding}. At the same time, networks tend to learn examples with consistent patterns before they memorise noisy ones~\cite{arpit2017closer}, and this ordering is the basis of many remedies: curricula and sample selection by small loss~\cite{jiang2017mentornet,han2018co,li2020dividemix}, regularisers that exploit early learning~\cite{liu2020early}, and data-independent regularisers such as mixup~\cite{zhang2017mixup} and label smoothing~\cite{mller2019when}. Comparisons run at larger scale and with per-method tuning make it hard to tell how much of a method's benefit comes from the method and how much from the recipe.

This paper is a deliberately small, controlled comparison of four training objectives on one tiny dataset: the scikit-learn 8$\times$8 digits. Using one fixed CNN, optimiser, schedule and data split, we compare plain cross-entropy (CE), label smoothing (LS), mixup, and a one-network small-loss selection rule in the spirit of co-teaching~\cite{han2018co}, at symmetric noise rates of 0, 20, 40 and 60\%. Besides clean test accuracy we measure memorisation directly on the corrupted training examples: the fraction predicted as their wrong label (\emph{mem\_rate}), the fraction predicted as their true label (\emph{recover\_rate}), and their difference (\emph{mem\_gap}).

We pre-registered five hypotheses (H1--H5, Sec.~\ref{sec:setup}) and report each as supported, refuted or inconclusive. We claim no new algorithm. The contributions are (i) an equal-budget, 5-seed comparison with paired tests, (ii) a direct memorised-versus-clean-label measure, and (iii) sensitivity sweeps for the hyperparameter of each method. All numbers come from CPU runs in the accompanying repository; the study is small and its limits are listed in Sec.~\ref{sec:limitations}.
